\chapter[Aproximación de Taylor]{Aproximación polinómica y fórmula de Taylor}
\label{chap:taylor}

La aproximación lineal sustituye una función por la recta que comparte su
valor y su primera derivada en un punto. Taylor prolonga esa idea: buscamos un
polinomio que comparta también derivadas de orden superior. El resultado es
finito y local. No debe confundirse con afirmar que la función coincide con
una serie infinita.

\section{Del contacto lineal al contacto de orden superior}
\label{sec:contacto-taylor}

\subsection{Orden de contacto}
\CanonicalUnit{CM-C06-U001}{}{}

Sean \(f\) y \(p\) funciones definidas cerca de \(a\). Decir que tienen
\emph{contacto de orden al menos \(n\)} en \(a\) significa, en este capítulo,
que
\[
  f^{(k)}(a)=p^{(k)}(a),
  \qquad k=0,1,\dots,n,
\]
si esas derivadas existen. El número \(n\) describe cuántas derivadas
coinciden; no es automáticamente el grado exacto de \(p\).

De forma equivalente, bajo la regularidad necesaria,
\[
  f(x)-p(x)=o\bigl(|x-a|^n\bigr)
  \qquad (x\to a).
\]
La notación \(o\) pequeña siempre debe declarar la variable, el centro y el
orden. Una afirmación de este tipo describe un comportamiento asintótico
local, no una igualdad exacta lejos de \(a\).

\subsection{Construcción y unicidad del polinomio}
\CanonicalUnit{CM-C06-U002}{}{}
\ProofRole{proof-028}{}

Supongamos que \(f\) es derivable \(n\) veces en \(a\). Buscamos un polinomio
escrito en potencias de \(x-a\),
\[
  p(x)=c_0+c_1(x-a)+\cdots+c_n(x-a)^n,
\]
que satisfaga \(p^{(k)}(a)=f^{(k)}(a)\) para \(k=0,\dots,n\). Al derivar y
evaluar en \(a\) obtenemos
\[
  p^{(k)}(a)=k!c_k,
\]
de modo que necesariamente
\[
  c_k=\frac{f^{(k)}(a)}{k!}.
\]

\begin{definicion}[Polinomio de Taylor]
El \emph{polinomio de Taylor de grado a lo sumo \(n\), centrado en \(a\)}, es
\[
  P_{n,a}(x)=\sum_{k=0}^{n}\frac{f^{(k)}(a)}{k!}(x-a)^k.
\]
Cuando \(a=0\), también se llama polinomio de Maclaurin.
\end{definicion}

La construcción prueba además la unicidad: ningún otro polinomio de grado a
lo sumo \(n\) puede compartir con \(f\) todas las derivadas hasta orden \(n\)
en \(a\). El grado efectivo puede ser menor si alguno de los coeficientes
finales se anula.

\subsection{Polinomios básicos}
\CanonicalUnit{CM-C06-U003}{}{}

Los polinomios de Maclaurin más utilizados son
\[
  e^x=\sum_{k=0}^{n}\frac{x^k}{k!}+R_n(x),
\]
\[
  \sin x=\sum_{k=0}^{m}(-1)^k\frac{x^{2k+1}}{(2k+1)!}
          +R_{2m+1}(x),
\]
\[
  \cos x=\sum_{k=0}^{m}(-1)^k\frac{x^{2k}}{(2k)!}
          +R_{2m}(x),
\]
y, para \(x>-1\),
\[
  \ln(1+x)=\sum_{k=1}^{n}(-1)^{k+1}\frac{x^k}{k}+R_n(x).
\]
Aquí cada suma es finita y \(R_n\) es el resto exacto correspondiente; estas
fórmulas no afirman todavía convergencia de una serie infinita.

El centro no tiene que ser cero. Por ejemplo, el polinomio de Taylor de
grado \(n\) de \(e^x\) centrado en \(a\) es
\[
  P_{n,a}(x)=e^a\sum_{k=0}^{n}\frac{(x-a)^k}{k!}.
\]

\section{La fórmula de Taylor y el resto}
\label{sec:formula-taylor-resto}

\subsection{Resto exacto y forma de Lagrange}
\CanonicalUnit{CM-C06-U004}{}{}
\ProofRole{proof-029}{}

Definimos el \emph{resto de Taylor} por
\[
  R_{n,a}(x)=f(x)-P_{n,a}(x).
\]
Es una diferencia exacta. La magnitud \(|R_{n,a}(x)|\), o una cota para ella,
se interpreta como error de aproximación.

\begin{teorema}[Taylor con resto de Lagrange]
Sea \(I\) un intervalo abierto, sean \(a,x\in I\), con \(a\ne x\), y sea
\(n\in\Nzero\). Supongamos que \(f\) y sus derivadas hasta orden \(n\) son
continuas en \(I\), y que \(f^{(n+1)}\) existe en el intervalo abierto cuyos
extremos son \(a\) y \(x\). Entonces existe un punto \(\xi\) estrictamente
entre \(a\) y \(x\)
tal que
\[
  f(x)=P_{n,a}(x)
       +\frac{f^{(n+1)}(\xi)}{(n+1)!}(x-a)^{n+1}.
\]
El último sumando es el \emph{resto de Lagrange}.
\end{teorema}

\begin{proof}[Estrategia de la demostración]
Se fija \(x\) y se introduce una función auxiliar que mide la diferencia
entre \(f\), su polinomio de Taylor en \(a\) y un múltiplo de
\((t-a)^{n+1}\). El múltiplo se elige para que la función auxiliar se anule
en \(t=a\) con sus primeras \(n\) derivadas y también en \(t=x\). Aplicando
Rolle de forma sucesiva se obtiene un punto \(\xi\) donde se anula su
derivada de orden \(n+1\). Al despejar el múltiplo resulta la fórmula
anterior. La idea esencial es transformar el resto en una igualdad accesible
mediante aplicaciones repetidas de Rolle.
\end{proof}

La regularidad se vincula al resultado concreto. No se supone, sin decirlo,
que una función derivable indefinidamente sea analítica ni que el resto tienda
a cero cuando \(n\to\infty\).

\subsection{Cotas del error}
\CanonicalUnit{CM-C06-U005}{}{}
\ProofRole{proof-030}{}

Si, además de las hipótesis del teorema, se conoce una cota
\[
  |f^{(n+1)}(t)|\le M
\]
para todo \(t\) entre \(a\) y \(x\), entonces
\[
  |R_{n,a}(x)|
  \le \frac{M}{(n+1)!}|x-a|^{n+1}.
\]
La cota debe cubrir todo el segmento relevante, no solo el centro.

Una aproximación de grado mayor no tiene por qué mejorar en cualquier punto
y bajo cualquier elección de cotas. Para comparar dos grados hay que estimar
sus restos en el mismo punto o en el mismo intervalo. La disminución del
término \(|x-a|^{n+1}/(n+1)!\) puede quedar compensada por el crecimiento de
las derivadas.

\subsection{Tres restos representativos}
\CanonicalUnit{CM-C06-U006}{}{}
\ProofRole{proof-030}{}

Para la exponencial centrada en cero,
\[
  R_n(x)=\frac{e^\xi}{(n+1)!}x^{n+1}
\]
para algún \(\xi\) entre \(0\) y \(x\). Para el seno y el coseno, todas las
derivadas tienen valor absoluto a lo sumo uno; por ello, con el polinomio de
Taylor del grado correspondiente,
\[
  |R_n(x)|\le\frac{|x|^{n+1}}{(n+1)!}.
\]
Para \(\ln(1+x)\), el intervalo usado para acotar las derivadas no puede
atravesar \(-1\), que es frontera del dominio.

\section{Aplicaciones controladas}
\label{sec:aplicaciones-taylor}

\subsection{Aproximación numérica con error garantizado}
\CanonicalUnit{CM-C06-U007}{}{}
\ProofRole{proof-030}{}

Para aproximar \(e=e^1\), tomemos el polinomio de Maclaurin de grado siete:
\[
  P_7(1)=\sum_{k=0}^{7}\frac1{k!}
        =\frac{685}{252}
        \approx 2{,}718254.
\]
Como \(e^t\le e\le3\) en \([0,1]\), donde la cota elemental \(e\le3\) se
obtiene de la sucesión definitoria mediante la expansión binomial y
\(k!\ge2^{k-1}\) para \(k\ge2\), el resto de Lagrange satisface
\[
  |e-P_7(1)|\le\frac{3}{8!}<7{,}5\cdot10^{-5}.
\]
La aproximación y la garantía son datos distintos: el decimal procede del
polinomio, mientras que el intervalo certificado procede de la cota del
resto.

\subsection{Una desigualdad local}
\CanonicalUnit{CM-C06-U008}{}{}
\ProofRole{proof-030}{}

Para \(|x|\le\pi\), el polinomio de Maclaurin de grado dos del coseno es
\(1-x^2/2\). La forma de Lagrange da
\[
  \cos x=1-\frac{x^2}{2}+\frac{\sin\xi}{6}x^3
\]
para algún \(\xi\) entre \(0\) y \(x\). Si \(x\ge0\), entonces
\(0\le\xi\le x\le\pi\), de modo que \(\sin\xi\ge0\). Si \(x\le0\),
entonces \(-\pi\le x\le\xi\le0\), por lo que \(\sin\xi\le0\) y
\(x^3\le0\). En ambos casos el resto es no negativo y
\[
  \cos x\ge1-\frac{x^2}{2}
  \qquad (|x|\le\pi).
\]

\subsection{Límites mediante desarrollos finitos}
\CanonicalUnit{CM-C06-U009}{}{}

Cuando aparecen cancelaciones, un desarrollo hasta el primer término no nulo
puede mostrar el orden dominante. Por ejemplo,
\[
  \sin x=x-\frac{x^3}{6}+o(x^3)
  \qquad(x\to0),
\]
y por tanto
\[
  \lim_{x\to0}\frac{x-\sin x}{x^3}=\frac16.
\]
El orden elegido debe superar todas las cancelaciones relevantes. Si basta
un límite notable, una equivalencia ya validada o l'Hôpital bajo sus
hipótesis, no es obligatorio usar Taylor; la herramienta se elige por la
estructura del problema.

\section{Álgebra de la notación \(o\) pequeña}
\label{sec:algebra-o-pequena}

\subsection{Suma, producto y composición}
\CanonicalUnit{CM-C06-U010}{}{}
\ProofRole{proof-032}{}

La relación \(r(x)=o((x-a)^n)\) cuando \(x\to a\) significa
\[
  \lim_{x\to a}\frac{r(x)}{(x-a)^n}=0,
\]
con la adaptación mediante valor absoluto cuando convenga. De la definición
se deducen, para órdenes compatibles,
\[
  o((x-a)^n)+o((x-a)^n)=o((x-a)^n)
\]
y
\[
  o((x-a)^m)\,o((x-a)^n)=o((x-a)^{m+n}).
\]

La sustitución en una composición necesita una regla explícita. Si
\(u(x)\to0\) y \(r(t)=o(t^n)\) cuando \(t\to0\), entonces
\[
  r(u(x))=o\bigl(u(x)^n\bigr)
\]
en los puntos donde la composición esté definida. En efecto,
\[
  \frac{r(u(x))}{u(x)^n}\longrightarrow0
\]
si \(u(x)\ne0\) cerca del punto; los ceros eventuales se tratan por la
extensión natural cuando \(r(0)=0\). Para reexpresar el resultado como una
potencia de \(x-a\) hace falta, además, conocer el orden de \(u(x)\). La
unicidad del polinomio de Taylor no justifica por sí sola esta sustitución.

\section{Enriquecimiento y límite de alcance}
\label{sec:restos-enriquecimiento}

\subsection{Resto de Cauchy}
\CanonicalUnit{CM-C06-U011}{}{}
\ProofRole{proof-031}{}

Como enriquecimiento, bajo las hipótesis adecuadas puede expresarse el resto
en la forma de Cauchy
\[
  R_{n,a}(x)
  =\frac{f^{(n+1)}(\xi)}{n!}(x-\xi)^n(x-a),
\]
donde \(\xi\) está entre \(a\) y \(x\). Esta forma no pertenece al núcleo de
aplicación del curso.

Existe también un resto integral, pero su enunciado y demostración dependen de
la integración, que no es un prerrequisito de este curso diferencial. Por
ello no se formula ni se usa aquí. Esta exclusión evita una dependencia
oculta y no debe interpretarse como una negación del resultado.

\subsection{Qué queda establecido}
\CanonicalUnit{CM-C06-U012}{}{}

El núcleo del capítulo es finito: construir \(P_{n,a}\), expresar el resto de
Lagrange, acotarlo y usarlo en aproximaciones, desigualdades y límites. No se
ha desarrollado una teoría de series de potencias ni se ha afirmado que toda
función suave coincida con su serie de Taylor.

Los ejercicios del material de partida y cualquier colección adicional
quedan fuera de esta fuente canónica y de la autorización P4.5B.
